% 1. Introducción
\begin{frame}{Introducción}
  \begin{itemize}
    \item SimAS 3.0 es un simulador educativo para comprender cómo funciona el análisis sintáctico descendente predictivo (LL(1)).
    \item Permite crear y editar gramáticas de contexto libre, construir la tabla predictiva y seguir, paso a paso, la derivación y el árbol sintáctico generado.
    \item La aplicación ha evolucionado desde SimAS 1.0 y 2.0, con una interfaz renovada, pestañas inteligentes, generación de informes y soporte multiidioma.
  \end{itemize}
\end{frame}

\begin{frame}{Motivación y alcance}
  \begin{itemize}
    \item Se partía de una versión previa con problemas de consistencia y portabilidad. Esta versión los corrige y unifica la experiencia para que no se pierda información al navegar ni al guardar.
    \item El trabajo se centra en el análisis descendente predictivo: eliminación de recursividad y factorización cuando procede, cálculo fiable de First/Follow y construcción de la tabla LL(1).
    \item Además, se incorporan funciones de recuperación de errores, informes completos y compatibilidad en Windows, macOS y Linux.
  \end{itemize}
\end{frame}

\begin{frame}{Qué aporta SimAS 3.0}
  \begin{itemize}
    \item Una forma visual y dinámica de aprender LL(1): la derivación y el árbol se construyen en tiempo real.
    \item Un asistente guiado que acompasa los pasos clave (transformaciones, conjuntos y tabla) con explicaciones y representación clara.
    \item Una base estable y coherente para la docencia: interfaz limpia, pestañas para trabajar con varias vistas y documentación exportable.
  \end{itemize}
\end{frame}


